\chapter{Listeria and listeriosis}
\index{Listeria and listeriosis}

\section*{Definition and Overview}
Listeriosis is a serious, foodborne bacterial infection caused by the pathogen \textit{Listeria monocytogenes}. Although listeriosis is relatively uncommon compared to other foodborne illnesses, it represents a major clinical and public health concern due to its high hospitalisation and mortality rates. In healthy individuals, the infection may cause only mild, transient gastrointestinal symptoms or remain completely asymptomatic. However, in vulnerable populations—specifically pregnant women, unborn foetuses, neonates, elderly individuals, and those with compromised immune systems—\textit{Listeria} can invade the bloodstream and central nervous system, leading to life-threatening complications such as septicaemia, meningitis, and chorioamnionitis.

\section*{Epidemiology}
Listeriosis is a relatively rare disease, with an annual incidence of approximately 0.5 to 1 case per 100,000 population in industrialised countries. In many countries, between 60 and 80 cases are notified annually. Despite its low incidence, listeriosis is characterised by an extremely high clinical severity; more than 90\% of diagnosed individuals require hospitalisation, and the mortality rate remains high, hovering between 20\% and 30\% globally. Outbreaks are typically associated with contaminated ready-to-eat foods, soft cheeses, deli meats, unpasteurised dairy products, pre-cut salads, and fresh fruits (such as rockmelon). The disease exhibits no strong seasonal pattern but occurs sporadically or in clusters linked to specific contaminated food production facilities.

\section*{Aetiology and Pathophysiology}
Listeriosis is caused by \textit{Listeria monocytogenes}, a Gram-positive, rod-shaped, facultative anaerobic bacterium. The organism's survival mechanisms and pathogenicity are unique:
\begin{itemize}
    \item \textbf{Environmental resilience:} \textit{Listeria monocytogenes} is exceptionally hardy. Unlike most other foodborne pathogens, it can survive and multiply at refrigeration temperatures (around 4\textdegree C), tolerate high salt concentrations, and survive in acidic environments. This allows it to persist and colonise food processing equipment.
    \item \textbf{Transmission:} the primary route of transmission is the consumption of contaminated foods. Vertical transmission from mother to foetus can occur transplacentally or during passage through the birth canal.
    \item \textbf{Intracellular invasion:} once ingested, the bacteria utilise specialised surface proteins called internalins (InlA and InlB) to bind to and invade the epithelial cells of the host's intestinal tract.
    \item \textbf{Phagosomal escape and actin-based motility:} after entering a host cell, the bacterium secretes a pore-forming toxin called listeriolysin O (LLO) and phospholipases to escape from the phagolysosome into the host cell cytoplasm. Once in the cytoplasm, it polymerises host cell actin via the bacterial surface protein ActA, generating "comet tails" that propel the bacterium directly into neighbouring cells.
    \item \textbf{Barrier crossing:} by moving directly from cell to cell, \textit{Listeria} evades humoral immune detection (antibodies) and successfully crosses the intestinal barrier, the blood-brain barrier, and the placental barrier.
\end{itemize}

\section*{Clinical Features}
The incubation period for listeriosis is highly variable, ranging from 1 to 70 days (typically 2 to 3 weeks). The clinical presentation depends heavily on the patient's immune status and risk profile:
\begin{itemize}
    \item \textbf{Healthy individuals:} may experience mild, self-limiting febrile gastroenteritis characterised by fever, muscle aches, headache, nausea, vomiting, and watery diarrhoea.
    \item \textbf{Pregnant individuals:} often present with a mild, flu-like illness consisting of fever, chills, fatigue, joint pain, and lower back pain. However, this mild maternal illness can be clinically catastrophic for the foetus, leading to spontaneous miscarriage, stillbirth, premature delivery, or severe neonatal infection.
    \item \textbf{Neonates (early-onset):} acquired in utero, presenting within hours of birth as sepsis, respiratory distress, and a characteristic disseminated granulomatous rash (granulomatosis infantiseptica), carrying a high mortality rate.
    \item \textbf{Neonates (late-onset):} acquired during delivery, presenting 1 to several weeks after birth, typically manifesting as bacterial meningitis or septicaemia.
    \item \textbf{Elderly and immunocompromised patients:} most commonly present with acute purulent meningitis or septicaemia. Symptoms include high fever, severe headache, neck stiffness, confusion, loss of balance, photophobia, tremors, and focal neurological deficits.
\end{itemize}

\section*{Differential Diagnosis}
Given its diverse presentations, listeriosis must be differentiated from other infections:
\begin{itemize}
    \item \textbf{Other bacterial meningitis:} infections caused by \textit{Streptococcus pneumoniae} or \textit{Neisseria meningitidis} present similarly, but \textit{Listeria} is more common in patients over 50, immunocompromised individuals, and pregnant patients.
    \item \textbf{Neonatal sepsis of other origins:} septicaemia caused by Group B Streptococcus (GBS), \textit{Escherichia coli}, or \textit{Streptococcus pneumoniae} must be differentiated through blood and cerebrospinal fluid cultures.
    \item \textbf{Influenza or viral syndrome in pregnancy:} a pregnant patient presenting with fever and muscle aches must be screened for listeriosis, as viral syndromes do not carry the same risk of bacterial transplacental transmission.
    \item \textbf{Other foodborne pathogens:} gastroenteritis from \textit{Salmonella}, \textit{Campylobacter}, or \textit{Shigella} is clinically similar but does not carry the high risk of systemic invasion and meningitis.
\end{itemize}

\section*{When to Seek Medical Care}
Pregnant women, elderly individuals, and immunocompromised patients who develop a fever, chills, muscle aches, or a severe headache—particularly if they have recently eaten high-risk foods or foods subject to a recall—should seek medical attention immediately.

\textbf{Contact your local emergency services for:}
\begin{itemize}
    \item Severe headache accompanied by neck stiffness, vomiting, and sensitivity to bright light
    \item Altered mental state, confusion, difficulty speaking, loss of coordination, or seizures
    \item High fever with rapid breathing, pale skin, shivering, or a rapid heart rate (signs of septicaemia)
    \item In a newborn: poor feeding, extreme lethargy, grunting or rapid breathing, or a bulging soft spot (fontanelle) on the head
\end{itemize}

\section*{Diagnosis and Investigations}
Listeriosis cannot be diagnosed on clinical grounds alone and requires microbiological confirmation:
\begin{itemize}
    \item \textbf{Blood cultures:} the primary diagnostic test for systemic listeriosis, which should be obtained in all patients presenting with fever and risk factors.
    \item \textbf{Cerebrospinal fluid (CSF) analysis:} performed via lumbar puncture in patients with suspected meningitis. CSF typically shows a high white cell count with neutrophilic or lymphocytic predominance, elevated protein, and low glucose. Gram stain may reveal Gram-positive rods, though they can be thin and easily mistaken for diphtheroids or streptococci.
    \item \textbf{PCR assays:} polymerase chain reaction tests on blood, CSF, or tissue samples offer rapid and highly sensitive detection of \textit{Listeria monocytogenes} DNA.
    \item \textbf{Culture of placental and foetal tissue:} in cases of miscarriage, stillbirth, or neonatal infection, cultures of the placenta, amniotic fluid, or foetal blood are critical.
    \item \textbf{Routine laboratory tests:} may reveal leucocytosis with a left shift, thrombocytopenia, and elevated inflammatory markers (such as C-reactive protein).
\end{itemize}

\section*{Management}
Treatment of listeriosis involves prompt, targeted antibiotic therapy and supportive care:
\begin{itemize}
    \item \textbf{First-line antibiotic therapy:} the treatment of choice is high-dose intravenous ampicillin or penicillin G, frequently combined with an aminoglycoside (such as gentamicin) for synergistic bactericidal activity.
    \item \textbf{Alternative therapy for penicillin-allergic patients:} trimethoprim-sulfamethoxazole (co-trimoxazole) is the primary alternative agent and crosses the blood-brain barrier effectively.
    \item \textbf{Duration of therapy:} treatment must be prolonged to prevent relapse. Systemic infections are treated for 2 to 3 weeks, while meningitis requires 3 to 4 weeks, and brain abscesses may require up to 6 weeks.
    \item \textbf{Supportive management:} includes intravenous fluids, electrolyte monitoring, management of intracranial pressure, and neonatal intensive care support for infected infants.
\end{itemize}

\section*{Complications}
Systemic \textit{Listeria} infection can lead to severe long-term sequelae or death, particularly in high-risk patients:
\begin{itemize}
    \item Miscarriage, stillbirth, or premature birth
    \item Neonatal death or severe congenital disabilities
    \item Permanent neurological damage in survivors of meningitis, including cranial nerve palsies, hydrocephalus, hearing loss, and developmental delay
    \item Septic shock and multi-organ failure
    \item Brain abscess or rhomboencephalitis (brainstem infection)
    \item Endocarditis, particularly in patients with pre-existing valvular heart disease
\end{itemize}

\section*{Prognosis and Outcomes}
The prognosis is highly dependent on the patient's age, comorbidities, and how quickly antibiotic therapy is initiated. In healthy adults, the prognosis is excellent, and gastroenteritis resolves completely without complications. In contrast, the prognosis is guarded for neonates and immunocompromised individuals, where the mortality rate remains around 20\% to 30\%. For pregnant women, while maternal recovery is almost always complete, the foetal prognosis is poor, with up to 50\% of cases resulting in miscarriage, stillbirth, or severe neonatal distress.

\section*{Prevention and Self-Care}
Vulnerable individuals can significantly reduce their risk of contracting listeriosis by adhering to strict food hygiene and dietary guidelines:
\begin{itemize}
    \item Avoid high-risk foods if you are pregnant, elderly, or immunocompromised. These include soft-serve ice cream, soft cheeses (such as brie, camembert, ricotta, and feta), pre-packaged deli meats, pâté, raw seafood (such as oysters and smoked salmon), pre-cut melons, and unpasteurised milk products
    \item Wash all fresh raw vegetables and fruits thoroughly under running water before peeling, cutting, or eating them
    \item Keep the refrigerator temperature set below 4\textdegree C and clean the inside of the refrigerator regularly to prevent bacterial colonisation
    \item Cook all raw meat, poultry, and seafood thoroughly to safe internal temperatures, and reheat leftovers until they are steaming hot
    \item Avoid cross-contamination by washing hands, knives, cutting boards, and countertops with warm, soapy water after preparing uncooked foods
    \item Consume perishable and ready-to-eat foods as soon as possible after opening, and discard any food that has exceeded its use-by date
\end{itemize}

\section*{Sources and Further Reading}
The following organisations provide reliable information on \textit{Listeria} and food safety:
\begin{itemize}
    \item Healthdirect Australia. \textit{Listeria infection and dietary precautions.} https://www.healthdirect.gov.au/
    \item Food Standards Australia New Zealand (FSANZ). \textit{Listeria food safety guidelines.} https://www.foodstandards.gov.au/
    \item Department of Health and Aged Care (Australia). \textit{Listeriosis clinical information and resources.} https://www.health.gov.au/
    \item Centers for Disease Control and Prevention. \textit{Listeria infection: symptoms, diagnosis, and prevention.} https://www.cdc.gov/
    \item World Health Organization. \textit{Listeriosis factsheet and food safety standards.} https://www.who.int/
\end{itemize}
